\section{Tại sao Deep Learning truyền thống chưa đủ?}

\textit{Tổng thời lượng: 3 phút 55 giây}


% ================= SCENE 1.1 =================
\subsection{Scene 1.1: Hook -- Giới hạn của DL hữu hạn chiều}
\textbf{Timestamp tổng:} \texttt{0:00 -- 0:45}

\noindent
\setlength{\tabcolsep}{4pt}
\renewcommand{\arraystretch}{1.2}
\begin{tabularx}{\textwidth}{|p{3.6cm}|>{\RaggedRight\arraybackslash}X|p{4.6cm}|}
\hline
\textbf{Visual (Manim)} & \textbf{Audio / Voiceover} & \textbf{Ghi chú Animation} \\ \hline

\textbf{[0:00--0:12]} \newline Màn hình đen $\to$ lưới $64 \times 64$ fade-in $\to$ từng ô sáng dần thành ảnh mèo. 
& \textit{"Trong hàng chục năm qua, Deep Learning đã tạo ra những đột phá không tưởng, thay đổi hoàn toàn diện mạo của thị giác máy tính và xử lý ngôn ngữ tự nhiên."} 
& Fade-in grid (0.5s) $\to$ pixel fill animation (1s/row) $\to$ image complete. Sync với nhịp VO. \\ \hline

\textbf{[0:12--0:28]} \newline Ảnh mèo morph thành tensor $\to$ ma trận trọng số $\to$ phép nhân ma trận trừu tượng. 
& \textit{"Chúng ta đã quá quen thuộc với việc CNN nhận diện hình ảnh, Transformer sinh văn bản, hay những kiến trúc kinh điển như ResNet và GPT."} 
& Morphing mượt: \manim{Transform} image $\to$ \manim{Matrix} $\to$ \manim{DotProduct}. Highlight keywords. \\ \hline

\textbf{[0:28--0:45]} \newline Zoom out: hiện khung $\mathbb{R}^n$ với các chấm điểm rời rạc. Overlay: \textbf{Finite-dimensional Euclidean Space}. 
& \textit{"Nhưng nếu lùi lại một bước để quan sát, tất cả các mô hình này đều chia sẻ một giới hạn bản chất: chúng được thiết kế để xử lý các đối tượng \textbf{hữu hạn chiều} trong không gian Euclid $\mathbb{R}^n$. Một bức ảnh, suy cho cùng, là một ma trận pixel hữu hạn. Một câu văn là một chuỗi các vector rời rạc."} 
& Camera pull-back (2s) $\to$ $\mathbb{R}^n$ frame \manim{Create} $\to$ text overlay \manim{Write}. Bold "hữu hạn chiều" sync với highlight. \\ \hline

\end{tabularx}

% ================= SCENE 1.2 =================
\subsection{Scene 1.2: Plot twist -- Bản chất dữ liệu: ``Đầu dò'' và ``Trường liên tục''}
\textbf{Timestamp tổng:} \texttt{0:45 -- 1:45}

\noindent
\setlength{\tabcolsep}{4pt}
\renewcommand{\arraystretch}{1.2}
\begin{tabularx}{\textwidth}{|p{3.6cm}|>{\RaggedRight\arraybackslash}X|p{4.6cm}|}
\hline
\textbf{Visual (Manim)} & \textbf{Audio / Voiceover} & \textbf{Ghi chú Animation} \\ \hline
\textbf{[0:45--1:05]}
\newline Khung $\mathbb{R}^n$ vỡ vụn (shatter). Từ các mảnh vỡ, một \textbf{Trường vô hướng 2D} (Scalar Field - vd: nhiệt độ) nổi lên dưới dạng một mặt cong 3D mượt mà (surface), uốn lượn liên tục.
\newline Một \textbf{"Đầu dò" (Probe)} phát sáng xuất hiện, bay tự do trong không gian. Đến đâu, nó thả xuống một tọa độ $(x,y)$ và nhổ lên một con số $T = 24.5^\circ C$.
& \textit{"Tuy nhiên, nếu bước ra khỏi Khoa học Máy tính, dữ liệu không nằm sẵn trong một ma trận pixel. Nó là một \textbf{trường liên tục} (continuous field) tồn tại ở khắp mọi nơi. Dù là nhiệt độ khí quyển, áp suất dưới lòng đất, hay vận tốc dòng khí... bản chất của chúng là một hàm số $u(x)$."}
& \manim{Shatter} $\mathbb{R}^n$ $\to$ \manim{Surface} morphs from shards.
\newline \textbf{Probe}: \manim{Dot} glowing, di chuyển theo \manim{Path}.
\newline \manim{DashedLine} từ Probe xuống mặt phẳng tọa độ, \manim{MathTex} $T(x,y)$ hiện ra và update realtime theo vị trí Probe. \\ \hline
\textbf{[1:05--1:25]}
\newline Chia đôi màn hình.
\newline \textbf{Trái (Camera/CNN):} Một lưới grid 64x64 (như một cái khuôn) ụp xuống trường liên tục. Nó chỉ "nhìn" thấy data qua các ô vuông cố định.
\newline \textbf{Phải (Neural Operator):} Không có lưới. Chỉ có hàm số $u(x)$ và hàng loạt các Probe bay vào truy vấn (query) ở bất kỳ tọa độ nào, kể cả những điểm nằm "giữa" các ô lưới cũ.
& \textit{"Một bức ảnh giống như một tấm lưới cứng được ụp xuống thế giới: bạn chỉ thấy những gì lọt vào ô lưới. Nhưng trong khoa học, chúng ta cần truy vấn (query) giá trị tại \textbf{bất kỳ tọa độ nào}, tính đạo hàm tại mọi điểm, và tích phân trên mọi miền. Dữ liệu không phải là ảnh. Dữ liệu là một \textbf{Toán tử} (Operator)."}
& \textbf{Left:} \manim{Grid} drops, locks onto surface. Các điểm ngoài grid bị mờ đi (fade opacity).
\newline \textbf{Right:} \manim{VGroup} của nhiều Probes bay loạn xạ, "chọc" vào mặt cong và lấy data. Overlay text: \textbf{Query Anywhere}. \\ \hline
\textbf{[1:25--1:45]}
\newline Hiện một \textbf{Nút vặn (Dial)} ghi chữ "Resolution".
\newline Vặn nút lên Max: 
\newline - Bên CNN: Lưới grid chia nhỏ ra chằng chịt $\to$ nặng nề, nhưng bản chất vẫn là các ô vuông rời rạc.
\newline - Bên NO: Mặt cong (hàm số) vẫn giữ nguyên shape, chỉ là ta đang "lấy mẫu" (sample) nó dày đặc hơn. Overlay: \textbf{Discretization Invariance}.
& \textit{"Và đây là điểm then chốt: khi ta cố ép một mô hình Deep Learning truyền thống xử lý dữ liệu này bằng cách rời rạc hóa... ta đã ném đi bản chất liên tục của nó. Neural Operator sinh ra để học chính cái \textbf{hàm số gốc}, chứ không phải học cái lưới rời rạc."}
& \manim{ValueTracker} điều khiển Dial.
\newline Grid bên trái \manim{Transform} thành lưới dày đặc (visual noise).
\newline Surface bên phải vẫn \manim{Smooth}, các Probe chỉ là xuất hiện nhiều hơn.
\newline Text \manim{Write} sync với nhịp "ném đi bản chất". \\ \hline
\end{tabularx}

% ================= SCENE 1.3 =================
\subsection{Scene 1.3: Grid Mismatch \& Ràng buộc vật lý}
\textbf{Timestamp tổng:} \texttt{1:45 -- 2:30}

\noindent
\setlength{\tabcolsep}{4pt}
\renewcommand{\arraystretch}{1.2}
\begin{tabularx}{\textwidth}{|p{3.6cm}|>{\RaggedRight\arraybackslash}X|p{4.6cm}|}
\hline
\textbf{Visual (Manim)} & \textbf{Audio / Voiceover} & \textbf{Ghi chú Animation} \\ \hline
\textbf{[1:45--2:00]} \newline Hai lưới cạnh nhau: $64 \times 64$ (thô, xanh) và $128 \times 128$ (mịn, tím), cùng sample 1 hàm nhiệt độ. & \textit{"Sự khác biệt này tạo ra một khoảng cách công nghệ khổng lồ. Nếu ta cố ép một mô hình Deep Learning truyền thống để dự báo thời tiết bằng cách rời rạc hóa khí quyển thành một lưới 64x64, ta sẽ lập tức vấp phải sai lầm."} & Side-by-side \manim{Grid} with color coding. Same function sampled differently: highlight mismatch in density. \\ \hline
\textbf{[2:00--2:15]} \newline Animation: model train trên lưới thô $\to$ đưa lưới mịn vào $\to$ output glitch/nhiễu + overlay đỏ \textbf{GRID MISMATCH}. & \textit{"Vấn đề là: nếu bạn huấn luyện mô hình trên lưới 64x64 nhưng lại muốn chạy thử nghiệm trên lưới mịn hơn như 128x128, mô hình sẽ thất bại hoàn toàn. Đây chính là hiện tượng \textit{grid mismatch}: mỗi lần thay đổi độ phân giải, bạn buộc phải huấn luyện lại từ đầu."} & Glitch effect: \manim{ValueTracker} noise on output. Red X \manim{Flash} + text \manim{Write}. Sync "grid mismatch" with overlay. \\ \hline
\textbf{[2:15--2:30]} \newline Hiện $\nabla u$ và $\int u \, dx$ với icon vật lý. Text overlay: \textbf{Discretization Dependence = Barrier to Science}. & \textit{"Hơn nữa, trong khoa học, các đạo hàm và tích phân của đầu ra phải tuân thủ nghiêm ngặt các định luật vật lý. Nhưng nếu đầu ra chỉ là một vector số rời rạc, làm sao chúng ta kiểm soát được tính nhất quán đó?"} & Equations \manim{Write} term-by-term. Physics icon \manim{Pulse}. Text overlay \manim{FadeIn} + subtle \manim{Pulse}. \\ \hline
\end{tabularx}


% ================= SCENE 1.4 =================
\subsection{Scene 1.4: Traditional Solvers \& Nỗi đau của Giải tích số}
\textbf{Timestamp tổng:} \texttt{2:30 -- 3:30}

\noindent
\setlength{\tabcolsep}{4pt}
\renewcommand{\arraystretch}{1.2}
\begin{tabularx}{\textwidth}{|p{3.6cm}|>{\RaggedRight\arraybackslash}X|p{4.6cm}|}
\hline
\textbf{Visual (Manim)} & \textbf{Audio / Voiceover} & \textbf{Ghi chú Animation} \\ \hline
\textbf{[2:30--2:50]} \newline Một trường vector liên tục (mượt mà) bị một tấm lưới (Grid) đè xuống. Các gradient bị "băm" thành các ô vuông rời rạc (Pixelation). Phương trình PDE hiện lên như một cái khuôn ép. & \textit{"Trước khi có Neural Operators, các nhà khoa học giải quyết những bài toán này bằng phương pháp số truyền thống. Bản chất là ép một thế giới liên tục vào một chiếc khuôn rời rạc: Phương trình Đạo hàm Riêng (PDE) được xấp xỉ bằng Sai phân hữu hạn trên một lưới cố định."} & \manim{ContinuousField} morphs thành \manim{DiscreteBlocks}. Grid lines \manim{Create} đè lên. Cảm giác "mất mát" chi tiết vật lý. \\ \hline
\textbf{[2:50--3:10]} \newline Animation "Solver Sweep": Một tia sáng quét qua lưới để giải lặp. Lưới nhân đôi độ mịn $\to$ tia sáng quét chậm rì. Biểu đồ Cost bên cạnh vẽ đường cong $O(N^3)$ vọt lên và đâm thủng trần màn hình. & \textit{"Vấn đề là, solver truyền thống phải giải lặp đi lặp lại. Khi độ phân giải tăng lên để bắt được các chi tiết vật lý nhỏ, chi phí tính toán không tăng tuyến tính, mà bùng nổ theo lũy thừa. Ta chạm phải một bức tường tính toán không thể vượt qua."} & \manim{ScanLine} di chuyển. Grid \manim{Transform} dày đặc. \manim{Axes} vẽ đường cong hàm mũ \manim{Create} và \manim{Flash} đỏ khi chạm đỉnh. \\ \hline
\textbf{[3:10--3:30]} \newline Hiện "Black Box" Solver. Mũi tên Forward (Xanh) đi qua. Mũi tên Gradient/Inverse (Đỏ) cố đi ngược lại nhưng đập vào bức tường "Non-differentiable" và vỡ tan. Overlay: \textbf{Bổ sung ML, không thay thế}. & \textit{"Hơn nữa, các solver này thường là những hộp đen không khả vi. Bạn không thể dễ dàng lấy đạo hàm để giải các bài toán ngược (inverse problems) hay tối ưu hóa thiết kế. Mục tiêu của ML không phải là vứt bỏ solver, mà là bổ sung một công cụ linh hoạt hơn vào toolbox."} & Mũi tên đỏ \manim{MoveAlongPath} $\to$ \manim{Shatter} khi chạm hộp đen. Text overlay \manim{Write} chậm, dứt khoát. \\ \hline
\end{tabularx}

% ================= SCENE 1.5 =================
\subsection{Scene 1.5: Câu hỏi then chốt \& Teaser Neural Operators}
\textbf{Timestamp tổng:} \texttt{3:30 -- 4:00}

\noindent
\setlength{\tabcolsep}{4pt}
\renewcommand{\arraystretch}{1.2}
\begin{tabularx}{\textwidth}{|p{3.6cm}|>{\RaggedRight\arraybackslash}X|p{4.6cm}|}
\hline
\textbf{Visual (Manim)} & \textbf{Audio / Voiceover} & \textbf{Ghi chú Animation} \\ \hline
\textbf{[3:30--3:45]} \newline Màn hình chia đôi: Trái $\mathbb{R}^n \to \mathbb{R}^m$ (vector), Phải $\mathcal{A} \to \mathcal{U}$ (hàm số). & \textit{"Câu hỏi đặt ra là: Liệu Machine Learning có thể học được trực tiếp ánh xạ từ không gian hàm này sang không gian hàm kia mà không cần giải PDE thủ công mỗi lần hay không?"} & Split screen: \manim{VGroup} with \manim{Transform} labels. Pause 0.5s after question. \\ \hline
\textbf{[3:45--3:55]} \newline Dấu hỏi lớn "nổ" thành logo Neural Operator (chữ N cách điệu + $\int$). & \textit{"Và câu trả lời chính là \textbf{Neural Operators} — thế hệ AI mới học trên không gian vô hạn chiều."} & Explosion: \manim{Flash} $\to$ logo \manim{Create} with subtle glow. Sync "Neural Operators" with reveal. \\ \hline
\textbf{[3:55--4:00]} \newline Fade to black + text: \textbf{Thế giới hàm số}. & \textit{(nhạc chuyển cảnh, pause ngắn)} & Clean \manim{FadeOut} $\to$ text \manim{Write} centered. Hold 1s before cut. \\ \hline
\end{tabularx}



\section{THẾ GIỚI HÀM SỐ}
\textit{Tổng thời lượng: 3 phút 30 giây}

\subsection{Scene 2.1: Bản chất dữ liệu khoa học -- Không gian, Hình học và Trường}
\textbf{Timestamp tổng:} \texttt{4:00 -- 5:00}

\noindent 
\setlength{\tabcolsep}{4pt}
\renewcommand{\arraystretch}{1.2}
\begin{tabularx}{\textwidth}{|p{3.6cm}|>{\RaggedRight\arraybackslash}X|p{4.6cm}|}
\hline
\textbf{Visual (Manim)} & \textbf{Audio / Voiceover} & \textbf{Ghi chú Animation} \\ \hline
\textbf{[4:00--4:20]} \newline 
Một khối khí quyển 3D (Fluid Volume) trong suốt. Bên trong là các đường dòng (streamlines) uốn lượn và gradient nhiệt độ. 
Một "Thấu kính Toán tử" (Operator Lens) quét qua, biến trạng thái $t$ thành $t+\Delta t$. 
Overlay: \textbf{Tiến hóa thời gian \& Bảo toàn}.
& \textit{"Hãy quên những bức ảnh đi. Trong khí tượng hay cơ học chất lưu, dữ liệu là một khối không gian 3D chứa các trường vector và vô hướng. Bài toán dự báo thực chất là một toán tử tiến hóa: ánh xạ toàn bộ trạng thái hôm nay sang ngày mai. Và output bắt buộc phải là một hàm liên tục, vì ta cần tính divergence hay curl để kiểm tra xem khối lượng và năng lượng có được bảo toàn hay không."} 
& \texttt{VGroup} các \texttt{Streamline} mượt mà. 
\texttt{Surface} gradient nhiệt độ mờ. 
Toán tử \texttt{Transform} block $t \to t+1$. 
Hiện biểu tượng $\nabla \cdot \mathbf{u} = 0$ (divergence-free) nhấp nháy nhẹ. \\ \hline
\textbf{[4:20--4:40]} \newline 
Chuyển sang lát cắt địa chất 2D/3D (Stratified layers). 
Một điểm nổ (earthquake) tạo ra các vòng tròn sóng (wavefronts) lan truyền, khúc xạ qua các lớp. 
Trên mặt đất, các sensor vẽ ra các đồ thị 1D (seismograms). 
Mũi tên \textbf{INVERSE} ngược từ đồ thị 1D $\to$ cấu trúc 3D ẩn dưới sâu.
& \textit{"Trong địa chấn, ta đối mặt với bài toán ngược. Sóng địa chấn lan truyền trong lòng đất tuân theo phương trình Helmholtz. Ta chỉ thu thập được các hàm tín hiệu 1D trên bề mặt. Thách thức là: từ không gian hàm 1D trên mặt, suy ngược ra không gian hàm 3D của cấu trúc địa chất ẩn bên dưới. Một ánh xạ giữa hai miền không gian hoàn toàn khác biệt."} 
& \texttt{Circle} wavefronts lan tỏa, \texttt{Bend} khi qua ranh giới lớp.
Sensor \texttt{Plot} ra hàm 1D dao động.
Mũi tên đỏ \texttt{GrowArrow} ngược từ $1D \to 3D$. Overlay: \textbf{Cross-Domain Mapping}. \\ \hline
\textbf{[4:40--5:00]} \newline 
Profile cánh máy bay (Airfoil) trừu tượng. 
Hình dáng vật thể (Geometry) là một hàm biên. 
Trường áp suất (Pressure field) bám sát bề mặt (Boundary layer) và tách dòng (Wake).
Chuyển cảnh: Hàm sóng xác suất (Schrödinger) tiến hóa liên tục.
Text: \textbf{Geometry $\to$ Field, Time $\to$ Continuum}.
& \textit{"Khí động học: hình dáng vật thể chính là một hàm số mô tả biên, và output là trường áp suất bám sát lớp biên nơi đạo hàm biến thiên cực mạnh. Động lực học phân tử: quỹ đạo không phải là các frame rời rạc, mà là một hàm liên tục trong không gian pha. Điểm chung: Input là Hàm, Output là Hàm, và Hình học của miền xác định mọi thứ."}
& \texttt{ParametricCurve} Airfoil. 
\texttt{Heatmap} áp suất rực rỡ ở lớp biên (nơi đạo hàm lớn).
Chuyển morph sang \texttt{ProbabilityCloud} của hàm sóng.
Chốt scene bằng text \texttt{Write} đồng bộ nhịp VO cuối. \\ \hline
\end{tabularx}

\subsection{Scene 2.2: Lưu ý then chốt -- Ảnh không phải là Hàm số \& Sự sụp đổ số}
\textbf{Timestamp tổng:} \texttt{5:00 -- 6:00} \textit{(Đã mở rộng từ 40s lên 60s để trình bày đầy đủ hơn)}


\noindent

\setlength{\tabcolsep}{4pt}
\renewcommand{\arraystretch}{1.2}
\begin{tabularx}{\textwidth}{|p{3.6cm}|>{\RaggedRight\arraybackslash}X|p{4.6cm}|}
\hline
\textbf{Visual (Manim)} & \textbf{Audio / Voiceover} & \textbf{Ghi chú Animation} \\ \hline
\textbf{[5:00--5:25]} \newline Hiện heatmap trường nhiệt độ (đỏ nóng, xanh lạnh). Zoom cực gần $\to$ lộ rõ các ô màu/pixel rời rạc. Gạch chéo chữ \textbf{ẢNH}, hiện chữ \textbf{HÀM SỐ LẤY MẪU}. 
& \textit{"Có một lưu ý quan trọng: Khi hình dung một hàm số như trường nhiệt độ, ta thường vẽ nó dưới dạng ảnh màu: chỗ đỏ là nóng, chỗ xanh là lạnh. Nhưng đó không phải là ảnh. Đó là hàm số được hiển thị, được lấy mẫu tại một số điểm nhất định để ta có thể nhìn thấy."} 
& Heatmap \manim{ImageMobject}. Zoom camera \manim{MoveTo}. Grid lines \manim{Create} over pixels. Strike-through effect on "ẢNH", \manim{Transform} to "HÀM SỐ LẤY MẪU". \\ \hline
\textbf{[5:25--5:40]} \newline Cảnh báo: Heatmap bị ép vào lưới 64x64 $\to$ hiện dấu chặn/red barrier. Overlay: \textbf{Ném đi bản chất liên tục $\to$ Bị ràng buộc lưới}. 
& \textit{"Sự khác biệt này cực kỳ quan trọng. Nếu ta xử lý nó như ảnh, đưa vào CNN và train trên lưới 64x64, là ta đang ném đi bản chất liên tục của dữ liệu. Và ta bị ràng buộc trong lưới đó."} 
& CNN icon/text \manim{FadeIn}. Red barrier \manim{Create} blocking output. Warning text \manim{Write} with subtle shake effect. \\ \hline
\end{tabularx}

\subsection{Scene 2.3: Từ Deep Learning truyền thống đến Operator Learning}
\textbf{Timestamp tổng:} \texttt{5:40 -- 6:40}

\noindent
\setlength{\tabcolsep}{4pt}
\renewcommand{\arraystretch}{1.2}
\begin{tabularx}{\textwidth}{|p{3.6cm}|>{\RaggedRight\arraybackslash}X|p{4.6cm}|}
\hline
\textbf{Visual (Manim)} & \textbf{Audio / Voiceover} & \textbf{Ghi chú Animation} \\ \hline
\textbf{[5:40--6:10]} \newline
\textbf{Split screen: Solver Loop vs Operator Forward}

\textbf{Trái (Solver):} Input $a_1$ $\to$ PDE box $\to$ iterative solver (vòng lặp Gauss-Seidel) $\to$ output $u_1$. Lặp lại cho $a_2$, $a_3$...

\textbf{Phải (Operator):} Một "Learned Operator G" lớn. Nhiều inputs $a_1, a_2, a_3$ $\to$ G $\to$ outputs $u_1, u_2, u_3$ (chỉ 1 forward pass mỗi cái).

Overlay: \textbf{Solve từng cái vs Học một lần}.
& \textit{"Deep Learning truyền thống học map từ vector sang vector. Nhưng trong scientific computing, ta đối mặt với một câu hỏi khác: thay vì giải PDE từ đầu cho mỗi input mới, liệu ta có thể học luôn solution operator? Solver truyền thống phải lặp lại quá trình tính toán đắt đỏ cho từng trường hợp. Operator Learning amortize chi phí đó: train một lần, inference nhanh hơn hàng nghìn lần."} 
& \textbf{Trái:} Solver box với \manim{ValueTracker} vòng lặp chạy. Input mới $\to$ solver chạy lại.

\textbf{Phải:} Operator G \manim{FadeIn}. Nhiều inputs bay vào, outputs bay ra mượt mà (chỉ 1 mũi tên mỗi cái).

\manim{Transform} solver loops $\to$ compressed thành operator. Text overlay \manim{Write} sync "amortize". \\ \hline
\textbf{[6:10--6:40]} \newline
\textbf{Darcy Flow: Gallery of Solutions}

Hiện một "gallery" gồm 6-8 cặp $(a_i, u_i)$ - các trường permeability và pressure tương ứng.

Animation: Từ gallery này, một mũi tên lớn "Operator G" xuất hiện, connecting tất cả.

Cuối cùng: Hiện phương trình Darcy mờ đi, overlay text: \textbf{Không giải PDE mỗi lần. Học ánh xạ G}.
& \textit{"Lấy ví dụ Darcy Flow: phương trình mô tả dòng chảy qua môi trường xốp. Input là hệ số khuếch tán $a(x)$, output là trường áp suất $u(x)$. Solver truyền thống phải giải PDE này cho mỗi $a$ mới. Operator Learning nhìn bài toán khác: ta có một tập các cặp $(a, u)$, và ta muốn học toán tử $G$ sao cho $G(a) \approx u$. Không giải PDE mỗi lần. Học ánh xạ."}
& Gallery \manim{FadeIn} từng cặp (a, u).

\manim{CurvedArrow} lớn "G" \manim{Create} connecting gallery.

Equation Darcy \manim{Write} mờ, rồi \manim{FadeOut}.

Text overlay \manim{Write}: "Học ánh xạ G". \\ \hline
\end{tabularx}


\subsection{Scene 2.4: Thách thức \& Định nghĩa Neural Operator}
\textbf{Timestamp tổng:} \texttt{6:40 -- 7:30}

\noindent

\setlength{\tabcolsep}{4pt}
\renewcommand{\arraystretch}{1.2}
\begin{tabularx}{\textwidth}{|p{3.6cm}|>{\RaggedRight\arraybackslash}X|p{4.6cm}|}
\hline
\textbf{Visual (Manim)} & \textbf{Audio / Voiceover} & \textbf{Ghi chú Animation} \\ \hline
\textbf{[6:40--7:00]} \newline Hiện 3 lưới khác nhau: vuông 64x64, tam giác bất quy tắc, mịn 128x128. Dấu chấm hỏi $\to$ hiện chữ \textbf{Discretization Invariance}. Tiếp theo: chấm query di chuyển tự do trong miền $\to$ \textbf{Continuous Query}. 
& \textit{"Việc học trong không gian hàm đòi hỏi giải quyết hai thách thức. Discretization Invariance: dữ liệu đo trên lưới khác nhau, model phải hoạt động trên mọi loại lưới không cần retrain. Continuous Query: đầu ra phải là hàm thực thụ, truy vấn được tại bất kỳ điểm nào để lấy đạo hàm kiểm tra vật lý, hoặc tính tích phân đo năng lượng."} 
& Grids \manim{Transform} sequentially. Question mark \manim{Flash} $\to$ text \manim{Write}. Query dot moves along \manim{Path}. Sync VO with each challenge reveal. \\ \hline
\textbf{[7:00--7:30]} \newline Hiện checklist 4 tiêu chí: 1) Input mọi độ phân giải. 2) Output truy vấn mọi điểm. 3) Hội tụ giới hạn liên tục. 4) Nhanh hơn solver hàng nghìn lần. 
& \textit{"Tóm lại, Neural Operator là framework thỏa mãn bốn tiêu chí: Nhận đầu vào là hàm số tại bất kỳ độ phân giải nào. Xuất đầu ra truy vấn được tại mọi điểm. Tính hội tụ: khi lưới càng mịn, mô hình hội tụ về giới hạn liên tục duy nhất. Và tốc độ: phải nhanh hơn hàng nghìn lần so với solver truyền thống. Phần tiếp theo: cấu trúc bên trong."} 
& Checklist items \manim{FadeIn} one by one (0.8s each). Subtle checkmark animation next to each item. Final hold 1s before section cut. \\ \hline
\end{tabularx}

\section{XÂY DỰNG NEURAL OPERATOR}
\textit{Tổng thời lượng: 3 phút 45 giây}

\subsection{Scene 3.1: Nền tảng Toán học -- Tích phân \& Đạo hàm}
\textbf{Timestamp tổng:} \texttt{7:30 -- 8:15}

\noindent
\setlength{\tabcolsep}{4pt}
\renewcommand{\arraystretch}{1.2}
\begin{tabularx}{\textwidth}{|p{3.6cm}|>{\RaggedRight\arraybackslash}X|p{4.6cm}|}
\hline
\textbf{Visual (Manim)} & \textbf{Audio / Voiceover} & \textbf{Ghi chú Animation} \\ \hline
\textbf{[7:30--7:55]} \newline Hiện đường cong $a(x)$ dưới dấu tích phân. Chia thành $N$ hình chữ nhật $\to$ $N$ tăng dần, $\Delta x \to 0$, tổng hội tụ thành diện tích mượt. 
& \textit{"Để hiểu Neural Operator, nhớ lại hai khái niệm cơ bản: Tích phân qua tổng Riemann. Chia hàm thành $N$ hình chữ nhật. Khi $N$ tiến tới vô cùng và $\Delta x$ tiến về 0, tổng các hình chữ nhật này sẽ hội tụ về giá trị tích phân thực."} 
& \manim{RiemannSum} animation. $N$ \manim{ValueTracker} tăng mượt. Area fill color transition từ rời rạc $\to$ liên tục. Equation \manim{Write} sync với VO. \\ \hline
\textbf{[7:55--8:15]} \newline Chuyển sang sai phân hữu hạn: 3 điểm trên đường cong $\to$ đường cát tuyến trượt thành tiếp tuyến. Hiện công thức đạo hàm rời rạc. 
& \textit{"Đạo hàm qua sai phân hữu hạn: Đạo hàm tại một điểm thực chất là sự chênh lệch giá trị giữa các điểm lưới chia cho khoảng cách."} 
& Secant line \manim{Transform} thành tangent. Points highlight sequentially. Formula appears term-by-term. Smooth easing. \\ \hline
\end{tabularx}

\subsection{Scene 3.2: Bước nhảy từ MLP sang Toán tử Tích phân}
\textbf{Timestamp tổng:} \texttt{8:15 -- 9:15}

\noindent
\setlength{\tabcolsep}{4pt}
\renewcommand{\arraystretch}{1.2}
\begin{tabularx}{\textwidth}{|p{3.6cm}|>{\RaggedRight\arraybackslash}X|p{4.6cm}|}
\hline
\textbf{Visual (Manim)} & \textbf{Audio / Voiceover} & \textbf{Ghi chú Animation} \\ \hline
\textbf{[8:15--8:45]} \newline 
Trên: Phương trình MLP. 
Dưới: Đồ thị rời rạc (các node nối nhau). 
Morphing: Các node tan chảy thành đường cong $a(x)$. Ma trận trọng số tan thành Kernel Heatmap $\kappa(y,x)$. Ký hiệu $\sum$ kéo giãn thành $\int$.
& \textit{"Hãy nhìn một layer MLP: một tổng có trọng số rời rạc. Bây giờ, hãy tưởng tượng số lượng node tiến tới vô cùng. Vector đầu vào tan chảy thành một hàm liên tục $a(x)$. Ma trận trọng số $W$ hòa tan thành một bề mặt Kernel liên tục $\kappa(y,x;\theta)$. Và phép tổng $\sum$ kéo giãn ra, trở thành tích phân $\int$."} 
& \manim{Graph} morphs thành \manim{ParametricCurve}. 
\manim{MatrixGrid} morphs thành \manim{Heatmap}. 
Ký hiệu $\sum$ \manim{Transform} mượt mà thành $\int$. \\ \hline
\textbf{[8:45--9:15]} \newline 
Hiện "Thấu kính Kernel" (Kernel Lens) quét dọc theo $a(x)$ để vẽ ra $v(y)$. 
Flash nhanh hình ảnh gợn sóng (Green's function). 
Chốt: \textbf{Kernel không còn cho trước, nó được học từ data}.
& \textit{"Về mặt hình học, toán tử tích phân giống như một thấu kính quét qua không gian: tại mỗi điểm $y$, nó gom thông tin từ toàn bộ hàm $a(x)$ thông qua kernel $\kappa$. Dạng toán này không mới — nó chính là hàm Green hay đáp ứng xung trong vật lý. Điểm khác biệt duy nhất: thay vì thiết kế kernel bằng tay, ta để mạng neural tự học nó từ dữ liệu."} 
& \manim{Window} quét dọc trục $x$, đồng thời \manim{Draw} đường cong $v(y)$. 
\manim{Flash} hình ảnh gợn sóng 2D. 
Text overlay \manim{Write} chậm, dứt khoát. \\ \hline
\end{tabularx}

\subsection{Scene 3.3: Từ điển Hình học \& Nghịch lý Làm mịn}
\textbf{Timestamp tổng:} \texttt{9:15 -- 10:15}

\noindent
\setlength{\tabcolsep}{4pt}
\renewcommand{\arraystretch}{1.2}
\begin{tabularx}{\textwidth}{|p{3.6cm}|>{\RaggedRight\arraybackslash}X|p{4.6cm}|}
\hline
\textbf{Visual (Manim)} & \textbf{Audio / Voiceover} & \textbf{Ghi chú Animation} \\ \hline
\textbf{[9:15--9:45]} \newline 
\textbf{Từ điển Hình học:} 
1. Cột chấm (Vector) mờ đi, Đường cong mượt (Hàm) mọc lên xuyên qua vị trí các chấm. 
2. Lưới ô vuông (Ma trận W) "tan chảy" (melt) thành Heatmap 2D mượt mà (Kernel $\kappa$). 
3. Ký hiệu $\sum$ kéo giãn morph thành $\int$.
& \textit{"Neural Operator không phải một kiến trúc kỳ lạ. Nó là sự dịch chuyển tự nhiên từ không gian hữu hạn sang vô hạn. Vector đầu vào tan chảy thành hàm số. Ma trận trọng số mở rộng thành một kernel liên tục. Và phép tổng rời rạc kéo giãn ra thành tích phân. Mọi thứ đều có một cặp song sinh trong không gian hàm."} 
& Không dùng text liệt kê. 
\manim{Transform} hoặc \manim{Morph} mượt mà. Nhấn mạnh sự tương đồng cấu trúc (structural isomorphism). \\ \hline
\textbf{[9:45--10:15]} \newline 
\textbf{Nghịch lý Làm mịn:} 
Hiện hàm Input có gai nhọn (Shock wave - nét đỏ gắt). 
Một \textbf{"Thấu kính Kernel" (Vầng hào quang vàng)} trượt dọc trục hoành. Khi lướt qua đỉnh nhọn, đỉnh nhọn bị \textbf{"bôi nhòe"} thành đường cong xanh mượt. 
Đột ngột, một \textbf{Tia Laser Residual} (Cyan rực rỡ) bắn thẳng từ đỉnh nhọn Input, xuyên thấu/vòng qua thấu kính, cắm thẳng vào Output. 
Kết quả: Đường cong mượt + Tia laser $\to$ \textbf{Gai nhọn phục hồi hoàn toàn}.
& \textit{"Nhưng tích phân mang theo một nghịch lý: nó là một bộ lọc thông thấp. Khi Kernel quét qua không gian, nó lấy trung bình có trọng số, vô tình xóa sổ các chi tiết sắc nét như sóng xung kích. Đây là lý do Residual Connection là yêu cầu sinh tồn. Nó bơm thẳng tín hiệu gốc bypass qua lõi tích phân, giữ lại những chi tiết vật lý quan trọng nhất."}
& \textbf{Kernel Lens:} \manim{Circle} mờ màu vàng trượt theo \manim{MoveAlongPath}. 
\textbf{Smoothing:} \manim{Transform} đường đỏ thành xanh khi Lens đi qua. 
\textbf{Residual:} \manim{ShowPassingFlash} hoặc \manim{Create} đường nét đứt Cyan rực rỡ. 
\textbf{Restoration:} \manim{Flash} tại điểm giao, khôi phục đường cong sắc nét. \\ \hline
\end{tabularx}

\subsection*{Scene 3.4: Deep Neural Operator \& Bẫy Rời rạc hóa}
\textbf{Timestamp tổng:} \texttt{10:15 -- 11:15}

\noindent
\setlength{\tabcolsep}{4pt}
\renewcommand{\arraystretch}{1.2}
\begin{tabularx}{\textwidth}{|p{3.6cm}|>{\raggedright\arraybackslash}X|p{4.6cm}|}
\hline
\textbf{Visual (Manim)} & \textbf{Audio / Voiceover} & \textbf{Ghi chú Animation} \\ \hline
\textbf{[10:15--10:45]} \newline 
\textbf{Pipeline Lift-Operate-Project:} \newline 
Input là 1 đường cong vật lý (xanh dương). \newline 
Encoder: Đường cong biến đổi thành một \textbf{dải lụa không gian đa chiều} rực rỡ (gradient xanh $\to$ tím). \newline 
Kernel Lens (vệt sáng vàng) quét ngang từ trái sang phải $\to$ dải lụa biến dạng uốn lượn theo. \newline 
Decoder: Dải lụa thu gọn lại thành 1 đường cong output duy nhất (xanh lá). \newline 
Overlay: \textbf{Dữ liệu luôn là Hàm số}.
& \textit{"Giống Deep Learning, ta xếp chồng nhiều lớp để tạo Deep Neural Operator. Nhưng kiến trúc đầy đủ tuân theo triết lý 'Lift - Operate - Project'. Encoder pointwise nâng hàm vật lý lên không gian đặc trưng ẩn. Các lớp toán tử tích phân trộn thông tin không gian mà không phá vỡ tính liên tục. Cuối cùng, Decoder chiếu về lại biến vật lý. Suốt quá trình, dữ liệu vẫn là hàm số, truy vấn được tại mọi điểm."} 
& \textbf{Lift:} \texttt{Create} dải lụa không gian (Surface/VMobject nhiều layer màu) từ 1 đường cong. \newline 
\textbf{Operate:} \texttt{Lens.animate.shift} quét ngang, dùng \texttt{Transform} biến đổi hình dạng dải lụa. \newline 
\textbf{Project:} \texttt{Transform} dải lụa $\to$ 1 đường output. \newline 
Thêm \texttt{Flash} hoặc các hạt \texttt{Particles} chạy dọc dải lụa. \\ \hline
\textbf{[10:45--11:15]} \newline 
\textbf{Bẫy Rời rạc hóa:} \newline 
Chia đôi màn hình. \newline 
Trên: Hàm trơn. $N$ tăng $\to$ chấm xanh dày đặc, error glow \textbf{co lại rất nhanh rồi biến mất}. \newline 
Dưới: Hàm Shock. $N$ tăng $\to$ chấm đỏ dày lên nhưng tại điểm gãy, error glow đỏ \textbf{co lại cực kỳ chậm}. \newline 
Hiện đồ thị hội tụ Log-Log: đường xanh dốc xuống rất nhanh, đường đỏ là một đường là là, độ dốc thấp. \newline 
Công thức $O(N^{-s})$.
& \textit{"Phân tích sai số của Neural Operator có ba thành phần: Xấp xỉ, Tổng quát hóa, và nhân tố mới: Rời rạc hóa. Sai số rời rạc hóa tuân theo $O(N^{-s})$, với $s$ là độ trơn Sobolev. Hàm càng trơn, sai số giảm càng nhanh khi lưới mịn dần. Nhưng với các nghiệm có shock wave hay gián đoạn, $s$ thấp, và sai số giảm rất chậm. Đây là lý do FNO tỏa sáng với dòng chảy trơn, nhưng cần các kiến trúc lai để bắt được các điểm gãy vật lý."} 
& \textbf{Sampling:} \texttt{Dot} + \texttt{DashedLine} nối với curve. \newline 
\textbf{Error glow:} \texttt{SurroundingRectangle} lồng trong loop tăng $N$, dùng \texttt{.animate.scale} giảm kích thước (giảm nhanh ở trên, giảm chậm ở dưới). \newline 
\textbf{Convergence plot:} \texttt{Axes} + \texttt{ax.plot} vẽ 2 đường thẳng có độ dốc (slope) khác nhau rõ rệt trên hệ trục Log-Log. \\ \hline
\end{tabularx}

\section{KIẾN TRÚC: FNO, GNO, U-NO}
\textit{Tổng thời lượng: 6 phút 10 giây}

\subsection{Scene 4.1: Fourier Neural Operator (FNO) -- Cơ sở tần số \& Định lý Tích chập}
\textbf{Timestamp tổng:} \texttt{11:15 -- 12:30}

\noindent
\setlength{\tabcolsep}{4pt}
\renewcommand{\arraystretch}{1.2}
\begin{tabularx}{\textwidth}{|>{\raggedright\arraybackslash}p{3.6cm}|>{\raggedright\arraybackslash}X|>{\raggedright\arraybackslash}p{4.6cm}|}
\hline
\textbf{Visual (Manim)} & \textbf{Audio / Voiceover} & \textbf{Ghi chú Animation} \\ \hline

\textbf{[11:15--11:50]} \newline 
\textbf{Phân rã sóng harmonic:} 
Hiện hàm tổng hợp $a(x) = \sin(x) + 0.5\sin(3x) + 0.3\sin(5x)$.
Tách thành 3 sóng riêng lẻ (mỗi sóng một màu) $\to$ chồng lên nhau.
Spectrogram thể hiện các mode tần số (1, 3, 5) với amplitude tương ứng.
& \textit{"Để tính tích phân hiệu quả, FNO không tính trực tiếp trên không gian thực. Thay vào đó, mô hình chuyển dịch hệ cơ sở sang miền tần số. Phép biến đổi Fourier phân rã hàm số thành tổ hợp sóng harmonic, cho phép ta quan sát dữ liệu dưới dạng phổ năng lượng với các hệ số phức. (ngắt 0.5s) Đây là bước đệm quan trọng nhất."} 
& \textbf{Phân rã sóng:} \manim{Transform} từ hàm tổng hợp $\to$ 3 sóng riêng lẻ. 
\textbf{Spectrogram:} \manim{BarChart} với height = amplitude, color = frequency. 
\textbf{Phổ năng lượng:} Highlight $\hat{a}(\omega)$ trong công thức FFT. 
\textbf{Chú ý:} Giữ hình 1.5s để viewer đọc công thức. \\ \hline

\textbf{[11:50--12:30]} \newline 
\textbf{Pipeline FNO \& Định lý Tích chập:} 
1. \textbf{FFT}: Hàm input $\to$ Spectrogram (màu xanh)
2. \textbf{Spectral Multiply}: Nhân với $R(\omega)$ (ma trận phức $d \times d$) $\to$ Highlight $R(\omega)$
3. \textbf{IFFT}: Spectrogram $\to$ Output hàm (màu xanh lá)
\textbf{So sánh:} 
- Tích chập trong miền thực: kernel di chuyển, $O(N^2)$
- Tích chập trong miền tần số: nhân phần tử, $O(N \log N)$
& \textit{"Trọng tâm tăng tốc nằm ở Định lý Tích chập: tích chập phức tạp trong không gian thực tương đương phép nhân phần tử trong miền tần số. Thay vì học kernel $O(N^2)$, ta học trực tiếp \textbf{ma trận trọng số $R(\omega)$} và nhân trong miền tần số. Một lớp FNO hoàn chỉnh gồm nhánh toàn cục học tương tác xa, và nhánh cục bộ $W a(x)$ bắt đặc trưng biến thiên nhanh. \textbf{Lưu ý quan trọng}: $R(\omega)$ là ma trận phức (không phải scalar) để học tương tác giữa các kênh. Nhờ FFT, độ phức tạp chỉ còn $O(N \log N)$."} 
& \textbf{Pipeline:} \manim{Create} tuần tự với \manim{Arrow} nối. 
\textbf{Highlight $R(\omega)$:} \manim{SurroundingRectangle} + text "Ma trận phức $d \times d$". 
\textbf{So sánh tốc độ:} Animation song song: 
- Trái: kernel quét, $O(N^2)$ 
- Phải: nhân phần tử, $O(N \log N)$ 
\textbf{Lưu ý:} Hiển thị $k_{max}$ truncation (chỉ giữ low-frequency modes). \\ \hline
\end{tabularx}

\subsection{Scene 4.2: Tính chất FNO \& Giới hạn}
\textbf{Timestamp tổng:} \texttt{12:30 -- 13:15}

\noindent
\setlength{\tabcolsep}{4pt}
\renewcommand{\arraystretch}{1.2}
\begin{tabularx}{\textwidth}{|p{3.6cm}|>{\RaggedRight\arraybackslash}X|p{4.6cm}|}
\hline
\textbf{Visual (Manim)} & \textbf{Audio / Voiceover} & \textbf{Ghi chú Animation} \\ \hline
\textbf{[12:30--13:15]} \newline 
Hiện 3 property cards: Tốc độ $O(N \log N)$, Tầm nhìn toàn cục (grid connected), Zero-shot resolution (64x64 $\to$ 256x256). 
Cuối cùng hiện warning: lưới đều/tuần hoàn + Gibbs phenomenon với shock wave.
& \textit{"FNO sở hữu ba tính chất đột phá. Thứ nhất: tốc độ xử lý lưới triệu điểm trong vài giây. Thứ hai: tầm nhìn toàn cục. Một phép nhân tần số kết nối mọi điểm cùng lúc, vượt xa CNN. Thứ ba: zero-shot generalization. Train 64x64, query 256x256 vẫn đạt độ chính xác chấp nhận được, vì FNO học ánh xạ hàm sang hàm. (ngắt 0.5s) Giới hạn: FNO phù hợp nhất với lưới đều và biên tuần hoàn. Nếu bài toán có shock wave, việc cắt giảm phổ tần số sẽ gây ra hiện tượng gợn sóng Gibbs, khiến sai số tăng vọt."} 
& Cards \manim{FadeIn} lần lượt, mỗi card giữ 0.8s. Global connectivity: \manim{Arrow} network lan tỏa. Resolution zoom \manim{Transform} mượt. Warning icon \manim{FadeIn} + pulse chậm. \\ \hline
\end{tabularx}

\subsection{Scene 4.3: Graph Neural Operator (GNO) -- Xử lý lưới bất quy tắc}
\textbf{Timestamp tổng:} \texttt{13:15 -- 14:30}

\noindent
\setlength{\tabcolsep}{4pt}
\renewcommand{\arraystretch}{1.2}
\begin{tabularx}{\textwidth}{|p{3.6cm}|>{\RaggedRight\arraybackslash}X|p{4.6cm}|}
\hline
\textbf{Visual (Manim)} & \textbf{Audio / Voiceover} & \textbf{Ghi chú Animation} \\ \hline
\textbf{[13:15--13:50]} \newline 
Mesh bất quy tắc. Xây đồ thị $G=(V,E)$ nối theo bán kính $r$. \newline
\textbf{Highlight 1 cell Voronoi} của node lân cận, tô sáng thể tích $w_j$. \newline
Text overlay: \textbf{GNN = Topo | GNO = Metric Space + Tích phân}. 
& \textit{"Thực tế kỹ thuật thường diễn ra trên miền bất quy tắc, nơi Fourier không áp dụng được. GNO chuyển dịch sang cấu trúc đồ thị. (ngắt 0.4s) Điểm khác biệt cốt lõi: GNN chỉ nhìn topo, còn GNO bảo toàn tính chất toán tử tích phân. Nó xấp xỉ tích phân bằng tổng có trọng số, trong đó $w_j$ chính là thể tích Voronoi của node nguồn. Khi lưới mịn dần, $w_j$ nhỏ lại, giúp mô hình hội tụ đúng về giới hạn liên tục mà không bị thiên lệch bởi mật độ điểm."} 
& Point cloud \manim{Create} $\to$ edges \manim{GrowFromCenter}. \newline
Dùng \manim{Fill} tô sáng 1 vùng Voronoi đại diện cho $w_j$. \newline
Text overlay \manim{Write} chậm. Nhấn mạnh sự đối lập Topo vs Metric. \\ \hline
\textbf{[13:50--14:30]} \newline Animation message passing: thông tin chạy dọc edges $\to$ tổng hợp tại node trung tâm $\to$ update. Hiện MLP kernel $\kappa_\theta$ box nhỏ trên edge. 
& \textit{"Kernel $\kappa$ không cố định mà là một MLP nhỏ tự học luật tương tác vật lý từ dữ liệu. GNO thực chất là Message Passing trong không gian hàm: node gửi thông điệp qua MLP, node nhận tổng hợp có trọng số Voronoi rồi áp dụng phi tuyến. Qua nhiều lớp, thông tin từ node xa nhảy cóc qua trung gian, giúp model nắm bắt tương tác toàn cục dù khởi đầu từ tổng hợp cục bộ."} 
& Data packets \manim{MoveAlongPath} tốc độ vừa phải. Aggregation \manim{Flash} tại center node. MLP box \manim{Scale} sync VO. Giữ frame cuối 0.8s trước khi cut. \\ \hline
\end{tabularx}


\subsection{Scene 4.4: U-shaped Neural Operator (U-NO) -- Đa tỷ lệ \& Skip Connection}
\textbf{Timestamp tổng:} \texttt{14:30 -- 15:45}

\noindent
\setlength{\tabcolsep}{4pt}
\renewcommand{\arraystretch}{1.2}
\begin{tabularx}{\textwidth}{|p{3.6cm}|>{\RaggedRight\arraybackslash}X|p{4.6cm}|}
\hline
\textbf{Visual (Manim)} & \textbf{Audio / Voiceover} & \textbf{Ghi chú Animation} \\ \hline
\textbf{[14:30--15:05]} \newline Kiến trúc chữ U: Encoder $\to$ Latent Space $\to$ Decoder. Highlight skip connections. Overlay hiệu ứng "làm mờ" của tích phân vs "giữ nét" của skip. 
& \textit{"Nghiệm PDE mang tính lưỡng cực: vừa có sóng dài toàn cục, vừa có gradient dốc đứng cục bộ như shock waves. FNO ưu tiên tần số thấp nên dễ làm mờ chi tiết. U-NO giải quyết mâu thuẫn này bằng kiến trúc phân cấp. Encoder thu nhỏ miền để nắm tương tác tầm xa với chi phí thấp. Latent Space nén trạng thái năng lượng tổng thể. Decoder khôi phục độ phân giải để tái tạo biến động sắc nét tại biên."} 
& U-structure \manim{Create} layer-by-layer. Downsampling \manim{Transform} kèm resolution label. Latent space \manim{Pulse} nhẹ. Skip arrows \manim{Create} màu sáng. Blur vs Sharp contrast effect chạy chậm 1.2s. \\ \hline
\textbf{[15:05--15:45]} \newline Minh họa skip connection: đặc trưng encoder bơm thẳng sang decoder. Hiện công thức $v_{dec} = \sigma(T_{dec} u_{dec} + u_{enc})$. Ẩn dụ vẽ bản đồ merge. 
& \textit{"Chìa khóa sinh tồn nằm ở Skip Connections. Bản chất của toán tử tích phân là một bộ lọc thông thấp — nó giỏi nắm bắt sóng dài nhưng lại vô tình 'bôi nhòe' các chi tiết sắc nét như shock waves hay lớp biên. *(Hiệu ứng đỉnh nhọn bị làm mờ)*. Các kết nối tắt này bơm thẳng đặc trưng hình học gốc từ encoder sang decoder, bypass qua lõi tích phân. *(Hiệu ứng đỉnh nhọn phục hồi)*. Công thức đảm bảo nghiệm cuối cùng vừa tuân thủ logic vật lý toàn cục, vừa giữ độ chính xác tuyệt đối tại điểm biến động. Giống vẽ bản đồ: cần toàn cảnh cho cao tốc, cần chi tiết cho tên đường nhỏ. U-NO làm cả hai cùng lúc."} 
& 1D Shock Wave \manim{Create} $\to$ \manim{Transform} thành Blurred func $\to$ Skip arrow \manim{GrowArrow} $\to$ Phục hồi sharp func. Formula \manim{Write} đồng bộ với text notes. Map metaphor visual transition giữ 1s. Cut nhẹ sang scene sau. \\ \hline
\end{tabularx}


\subsection{Scene 4.5: CoDANO, Spherical FNO \& Khối Đạo hàm}
\textbf{Timestamp tổng:} \texttt{15:45 -- 17:00}

\noindent
\setlength{\tabcolsep}{4pt}
\renewcommand{\arraystretch}{1.2}
\begin{tabularx}{\textwidth}{|p{3.6cm}|>{\RaggedRight\arraybackslash}X|p{4.6cm}|}
\hline
\textbf{Visual (Manim)} & \textbf{Audio / Voiceover} & \textbf{Ghi chú Animation} \\ \hline
\textbf{[15:45--16:20]} \newline Ma trận attention $\alpha_{c,c'}$ giữa các kênh vật lý (áp suất $\to$ vận tốc). Chuyển sang mặt cầu với lưới spherical harmonics thay sin/cos. 
& \textit{"Khi bài toán mở rộng lên hàng trăm biến tương tác, ta cần CoDANO: áp dụng cơ chế Attention giữa các biến vật lý, thay vì điểm không gian. Trọng số $\alpha_{c,c'}$ tự học định luật tương tác, ví dụ áp suất ảnh hưởng bao nhiêu phần trăm lên vận tốc. (ngắt 0.4s) Với khí hậu toàn cầu, Spherical FNO thay sin/cos bằng hàm điều hòa cầu, cho phép tích chập trực tiếp trên mặt cầu, triệt tiêu sai số tại hai cực."} 
& Attention weights \manim{Matrix} animate row/column chậm. Sphere \manim{ParametricSurface} với harmonic grid overlay. Camera orbit nhẹ 360° trong 2s. Giữ frame 0.5s. \\ \hline
\textbf{[16:20--17:00]} \newline Nghịch lý: tích phân làm mịn $\to$ mất đạo hàm. Giải pháp: khối CNN stencil $[-1, 0, 1]$ trượt trên hàm, tính sai phân hữu hạn. 
& \textit{"Tồn tại một nghịch lý: tích phân giỏi tổng hợp nhưng kém nắm bắt đạo hàm tức thời. Để tính lực cản hay dòng chảy, ta buộc phải tích hợp khối CNN. Bản chất bộ lọc CNN chính là phép sai phân hữu hạn tìm đạo hàm. Khi tăng độ phân giải, ta chỉ cần điều chỉnh hệ số stencil để các định luật vật lý luôn nhất quán. Sự kết hợp này bù đắp điểm mù của toán tử tích phân."} 
& Smooth curve vs derivative spike comparison hiện song song. CNN stencil \manim{Slide} dọc hàm, tốc độ vừa. Coefficient scaling \manim{Transform} khi resolution đổi. Hold 0.8s trước khi fade. \\ \hline
\end{tabularx}

\subsection{Scene 4.6: Tổng kết kiến trúc \& Chuyển cảnh}
\textbf{Timestamp tổng:} \texttt{17:00 -- 17:25}

\noindent
\setlength{\tabcolsep}{4pt}
\renewcommand{\arraystretch}{1.2}
\begin{tabularx}{\textwidth}{|p{3.6cm}|>{\RaggedRight\arraybackslash}X|p{4.6cm}|}
\hline
\textbf{Visual (Manim)} & \textbf{Audio / Voiceover} & \textbf{Ghi chú Animation} \\ \hline
\textbf{[17:00--17:25]} \newline Decision grid: FNO (lưới đều/tốc độ), GNO (mesh phức tạp), U-NO (đa tỷ lệ), CoDANO (nhiều biến). Badge \textbf{Hybrid SOTA}. List 4 thành phần: Tích phân + Residual + Bias hàm + CNN đạo hàm. Fade to black + text: \textbf{Section 5: Ứng dụng thực tế}. 
& \textit{"Tổng kết chiến lược: lưới đều cần tốc độ chọn FNO. Mesh phức tạp chọn GNO. Nghiệm đa tỷ lệ chọn U-NO. Nhiều biến tương tác chọn CoDANO. Thực tế SOTA luôn là hybrid. Một Neural Operator tối ưu thường phối hợp bốn thành phần: tích phân bắt tương tác xa, residual ổn định training, bias hàm số, và khối CNN xử lý đạo hàm. (ngắt 0.5s) Phần tiếp theo: những kiến trúc này đang thay đổi thế giới thực tế như thế nào."} 
& Grid cells \manim{FadeIn} kèm checkmark. Hybrid badge \manim{Create} glow nhẹ. 4 component icons \manim{Pulse} tuần tự. Clean \manim{FadeOut} $\to$ section title \manim{Write}. Hold 1.2s trước cut. \\ \hline
\end{tabularx}

\section{SECTION 5: ỨNG DỤNG THỰC TẾ}
\textit{Tổng thời lượng: 6 phút 45 giây}

\subsection{Scene 5.1: FourCastNet \& Spherical FNO -- Dự báo thời tiết}
\textbf{Timestamp tổng:} \texttt{17:25 -- 19:10}

\noindent
\setlength{\tabcolsep}{4pt}
\renewcommand{\arraystretch}{1.2}
\begin{tabularx}{\textwidth}{|p{3.6cm}|>{\RaggedRight\arraybackslash}X|p{4.6cm}|}
\hline
\textbf{Visual (Manim)} & \textbf{Audio / Voiceover} & \textbf{Ghi chú Animation} \\ \hline
\textbf{[17:25--18:10]} \newline Hiện cầu Trái Đất với các lớp dữ liệu thời tiết (nhiệt, gió). Mũi tên thời gian Hôm nay $\to$ Ngày mai. \newline \textbf{Cuộc đua tốc độ}: Thanh tiến trình đỏ (IFS) nhích cực chậm (1 Giờ) vs thanh xanh (FourCastNet) vọt lẹ tức thì (< 1 Giây). Bảng metrics (45,000x).
& \textit{"Bây giờ, hãy cùng xem Neural Operators thực sự thay đổi thế giới như thế nào. FourCastNet dựa trên FNO, train trên 40 năm dữ liệu ERA5. Nhiệm vụ: dự báo trạng thái khí quyển ngày mai từ hôm nay. So với mô hình IFS truyền thống: nhanh hơn 45,000 lần, tiết kiệm năng lượng 25,000 lần, độ chính xác tương đương ở dự báo ngắn và trung hạn."} 
& \manim{Circle} globe \& arcs. \textbf{Speed race}: \manim{Rectangle} \texttt{stretch\_to\_fit\_width}. IFS chậm (5s), FourCastNet chớp mắt (0.5s) + \manim{Flash} xanh. Metrics \manim{FadeIn}. \\ \hline
\textbf{[18:10--18:40]} \newline \textbf{Ensemble forecasting}: Đồ thị Nhiệt độ/Thời gian. Đường dự báo Deterministic (trắng) $\to$ bùng nổ thành đám mây 60 đường cong ngẫu nhiên (Spaghetti plot, Stochastic). Hiện dải xác suất upper/lower bounds.
& \textit{"Con số này có hệ quả thực tế: thay vì chạy một dự báo duy nhất, ta chạy hàng trăm kịch bản song song. Thay vì ngày mai 25 độ, model nói: 70 phần trăm khả năng 24 đến 26 độ, 20 phần trăm 22 đến 24 độ. Đây là thứ hệ thống cảnh báo thiên tai thực sự cần."} 
& \manim{Axes} plot. \manim{TransformFromCopy} đường trắng thành cụm \manim{ParametricFunction} opacity thấp (vàng). \manim{Create} dải xác suất đỏ/xanh. \\ \hline
\textbf{[18:40--19:10]} \newline \textbf{Khắc phục Singularity}: \newline Trái (Planar FNO): Lưới bị nhiễu/rách tại cực Bắc (Dấu X đỏ). \newline Phải (SFNO): Cơn bão di chuyển mượt mà qua cực. \newline Morphing phương trình Fourier $\to$ Spectral ($Y_l^m$).
& \textit{"Biến thể nâng cao là Spherical FNO. Khi dùng lưới phẳng chiếu lên cầu, model thường bị sai số kỳ dị (singularity) ở các cực. SFNO thay sin/cos bằng hàm điều hòa cầu. Cơn bão đi qua cực hoàn toàn mượt mà. Đáng chú ý: đây chính xác là cơ sở toán học domain experts dùng trong solver truyền thống. Neural Operator nói cùng ngôn ngữ với họ."} 
& \textbf{Planar}: Grid hội tụ + \manim{Cross} glitch. \textbf{SFNO}: \manim{MoveAlongPath} (cơn bão qua cực). \textbf{Morphing}: \manim{TransformMatchingTex} từ $\mathcal{G}_\theta$ sang $\sum \hat{u}_{l,m} Y_l^m$. \\ \hline
\end{tabularx}

\subsection{Scene 5.2: Carbon Storage \& Khí động học công nghiệp}
\textbf{Timestamp tổng:} \texttt{19:10 -- 20:55}

\noindent
\setlength{\tabcolsep}{4pt}
\renewcommand{\arraystretch}{1.2}
\begin{tabularx}{\textwidth}{|p{3.6cm}|>{\RaggedRight\arraybackslash}X|p{4.6cm}|}
\hline
\textbf{Visual (Manim)} & \textbf{Audio / Voiceover} & \textbf{Ghi chú Animation} \\ \hline
\textbf{[19:10--19:55]} \newline Hiện mô hình 3D vỉa đá dưới lòng đất. Dòng $CO_2$ bơm xuống $\to$ đám mây lan rộng theo thời gian. Hiện con số 700,000x. 
& \textit{"Một công cụ quan trọng chống biến đổi khí hậu là carbon capture: bơm $CO_2$ xuống vỉa đá sâu. Câu hỏi then chốt: bao nhiêu $CO_2$ lưu trữ an toàn, và đám mây sẽ lan đến đâu? Input là cấu trúc vận tốc 3D phức tạp. Output là sự lan rộng theo thời gian. Neural Operator nhanh hơn 700,000 lần so với solver truyền thống."} 
& 3D subsurface volume \manim{Create}. $CO_2$ plume \manim{Expand} with time counter. Speed metric \manim{Flash}. Smooth shading for depth. \\ \hline
\textbf{[19:55--20:25]} \newline Hiện hàng nghìn kịch bản lòng đất chồng mờ $\to$ bản đồ rủi ro tổng hợp. Chuyển sang geometry xe/cánh máy bay. 
& \textit{"Với tốc độ đó, ta chạy hàng nghìn kịch bản cấu trúc lòng đất khác nhau vì ta không biết chính xác lòng đất trông như thế nào. Điều này cho phép đánh giá rủi ro đầy đủ và đưa ra quyết định tốt hơn về vị trí lưu trữ."} 
& Multiple scenarios \manim{Overlap} then merge into risk heatmap. Quick pan to automotive/aerospace geometry. \\ \hline
\textbf{[20:25--20:55]} \newline Pipeline hybrid: Mesh 3D không đều $\to$ GNO $\to$ lưới đều $\to$ FNO (FFT) $\to$ decode ngược. Minh họa real-time design iteration. 
& \textit{"Đây là killer app công nghiệp. Quy trình cũ: thay đổi hình dáng $\to$ xây lưới mới $\to$ giải Navier-Stokes $\to$ mất nhiều giờ. Với Neural Operator: train model duy nhất, inference vài giây. Nhanh hơn 140,000 lần. Kiến trúc hybrid: GNO xử lý mesh không đều, map sang lưới đều, áp dụng FNO, decode ngược lại. Tối ưu thiết kế real-time."} 
& Pipeline blocks \manim{Create} left-to-right. Mesh morphing animation. Iteration counter speeds up. Clean hold. \\ \hline
\end{tabularx}

\subsection{Scene 5.3: Động lực học phân tử \& PINO}
\textbf{Timestamp tổng:} \texttt{20:55 -- 22:40}

\noindent
\setlength{\tabcolsep}{4pt}
\renewcommand{\arraystretch}{1.2}
\begin{tabularx}{\textwidth}{|p{3.6cm}|>{\RaggedRight\arraybackslash}X|p{4.6cm}|}
\hline
\textbf{Visual (Manim)} & \textbf{Audio / Voiceover} & \textbf{Ghi chú Animation} \\ \hline
\textbf{[20:55--21:40]} \newline So sánh: ML truyền thống (chấm rời $t \to t+\Delta t$) vs NO (đường cong liên tục). Hiện các bước trung gian của solver được dùng làm supervision. 
& \textit{"ML truyền thống xử lý động lực học phân tử theo kiểu rời rạc: state tại $t$ dự đoán state tại $t+\Delta t$. Nhưng bản chất bài toán là liên tục theo phương trình Schrödinger. Neural Operator học toán tử tiến hóa liên tục. Lợi ích lớn: dùng tất cả bước trung gian của solver làm supervision. Output là emulator domain expert kiểm tra được bằng cách cắm vào phương trình gốc."} 
& Discrete dots vs continuous curve comparison. Solver steps highlight sequentially. Equation constraint overlay. \\ \hline
\textbf{[21:40--22:10]} \newline Hiện cấu trúc PINO: $u_\theta$ cắm vào PDE $\to$ tính residual. Hiện công thức loss chia đôi: Data Loss + Physics Loss. 
& \textit{"Cho đến nay, ứng dụng đều học từ cặp data. Nhưng trong khoa học, ta còn biết định luật vật lý. PINO khai thác điều này: nếu model dự đoán $u_\theta$, cắm vào PDE kiểm tra thỏa mãn không. Nếu không, đó là sai số dùng để train. Loss gồm data loss supervised và physics loss tính bằng automatic differentiation."} 
& PDE residual check animation. Loss formula splits into two colored components. \manim{Write} syncs with VO. \\ \hline
\textbf{[22:10--22:40]} \newline Minh họa zero-shot upscaling: 64x64 $\to$ 256x256 chỉ bằng physics fine-tune, không cần data mới. Đường spectrum khớp ground truth. 
& \textit{"Kết quả ấn tượng: train FNO trên Darcy flow 64x64, fine-tune bằng physics loss không có thêm data. Kết quả 256x256 khớp chính xác ground truth dù model chưa bao giờ thấy data độ phân giải đó."} 
& Resolution zoom \manim{Transform}. Spectrum lines merge perfectly. Subtle glow on match point. \\ \hline
\end{tabularx}


\subsection{Scene 5.4: Mô hình sinh GANO \& Tổng kết}
\textbf{Timestamp tổng:} \texttt{22:40 -- 24:10}

\noindent
\setlength{\tabcolsep}{4pt}
\renewcommand{\arraystretch}{1.2}
\begin{tabularx}{\textwidth}{|p{3.6cm}|>{\RaggedRight\arraybackslash}X|p{4.6cm}|}
\hline
\textbf{Visual (Manim)} & \textbf{Audio / Voiceover} & \textbf{Ghi chú Animation} \\ \hline
\textbf{[22:40--23:25]} \newline Hiện pipeline GAN: Gaussian Process noise $\to$ Generator (NO) $\to$ hàm output. Discriminator nhận hàm $\to$ scalar. Nhiều sample tạo cloud phân phối + confidence band. 
& \textit{"Tất cả ứng dụng trên đều deterministic: cho input ra một output. Nhưng khoa học thường cần phân phối output. GANO tổng quát hóa GAN lên không gian hàm. Generator là NO nhận Gaussian Process sinh hàm đầu ra. Discriminator là Neural Functional nhận hàm trả về số thực. Khi hội tụ, generator học phân phối thật. Lấy mẫu nhiều lần $\to$ có phân phối dự đoán $\to$ khoảng tin cậy, variance, risk metrics. Hữu ích cho data augmentation, bài toán ngược Bayesian, khám phá không gian thiết kế."} 
& Noise paths \manim{Create} $\to$ Generator box $\to$ output functions stack. Distribution cloud forms. Confidence band shades. Smooth easing. \\ \hline
\textbf{[23:25--24:10]} \newline Tổng kết 3 ứng dụng chính trên màn hình: Weather, Carbon/CFD, Molecular/PINO. Fade sang text: \textbf{Section 6: Bài toán mở \& Tương lai}. 
& \textit{"Tóm lại: Neural Operators không chỉ là lý thuyết. Chúng đang thay đổi dự báo thời tiết, tối ưu carbon capture, tăng tốc thiết kế khí động học, mô phỏng phân tử liên tục, tích hợp vật lý qua PINO, và mô hình hóa bất định qua GANO. Phần tiếp theo: những bài toán mở và tương lai của AI for Science."} 
& Application icons/labels \manim{FadeIn} in grid. Clean \manim{FadeOut} $\to$ section title \manim{Write}. Hold 1s before cut. \\ \hline
\end{tabularx}

\section{NHỮNG BÀI TOÁN MỞ VÀ TƯƠNG LAI CỦA AI FOR SCIENCE}
\textit{Tổng thời lượng: 5 phút 15 giây}

\subsection{Scene 6.1: Open Problems -- Thiết kế kiến trúc \& Tích hợp vật lý}
\textbf{Timestamp tổng:} \texttt{24:10 -- 25:25}

\noindent
\setlength{\tabcolsep}{4pt}
\renewcommand{\arraystretch}{1.2}
\begin{tabularx}{\textwidth}{|p{3.6cm}|>{\RaggedRight\arraybackslash}X|p{4.6cm}|}
\hline
\textbf{Visual (Manim)} & \textbf{Audio / Voiceover} & \textbf{Ghi chú Animation} \\ \hline
\textbf{[24:10--24:45]} \newline
\textbf{Beat 1 -- Architecture Open Questions:} Hiện title ``Bài toán mở: Kiến trúc \& Vật lý''. Bên trái: ``Known Architectures'' với 4 ô nhỏ (FNO xanh, GNO cam, U-NO tím, CoDANO vàng). Bên phải: ``Open Questions'' với 4 ô dấu hỏi mờ dần hiện text (Max-pooling hàm? / Attention biến thể? / Scale 3D 10M+? / Convolution hàm?).
& \textit{``Chúng ta đã đi qua hầu như tất cả những gì đã biết về Neural Operators. Và điều thực sự hấp dẫn là: hầu như mọi thứ trong lĩnh vực này vẫn còn mở...''}
& Title \manim{FadeIn}. Known panel: 4 ô \manim{FadeIn} + \manim{Chip} style. Open panel: question marks \manim{Flash} $\to$ text \manim{Write} tuần tự. Sử dụng \manim{RoundedRectangle} cho mỗi ô, text \manim{move\_to} trong ô. \\ \hline
\textbf{[24:45--25:25]} \newline
\textbf{Beat 2 -- Physics Integration:} Clear toàn bộ beat 1. Hiện ``Tích hợp Vật lý: PINO \& Beyond''. Trái: PINO Pipeline (Input hàm $\to$ Neural Operator $\to$ Output $u_\theta$ $\to$ PDE Residual). Phải: 3 câu hỏi mở xếp dọc (Hard-code conservation? / Balance losses? / High-order derivatives?). Dưới cùng: Thanh tiến trình ``Potential vs Reality'' cho thấy gap.
& \textit{``Vấn đề mở thứ hai: tích hợp vật lý đúng cách...''}
& PINO pipeline \manim{Create} left-to-right. Question cards \manim{FadeIn} tuần tự. Progress bar \manim{GrowFromEdge}. Fade to black kết thúc. \\ \hline
\end{tabularx}

\subsection{Scene 6.2: Lý thuyết nền tảng \& Lượng hóa bất định}
\textbf{Timestamp tổng:} \texttt{25:25 -- 26:40}

\setlength{\tabcolsep}{4pt}
\renewcommand{\arraystretch}{1.2}
\begin{tabularx}{\textwidth}{|p{3.6cm}|>{\RaggedRight\arraybackslash}X|p{4.6cm}|}
\hline
\textbf{Visual (Manim)} & \textbf{Audio / Voiceover} & \textbf{Ghi chú Animation} \\ \hline
\textbf{[25:25--26:05]} \newline
\textbf{Beat 1 -- Knowledge Constellation:} Chia đôi màn hình. Trái ``Deep Learning Theory'': đồ thị dày đặc các node liên kết (Approximation, VC Dim, Rademacher, PAC-Bayes, NTK). Phải ``Neural Operator Theory'': vùng sương mù, 1 node sáng mờ (Universal Approximation) + 3 ``beacon'' nhấp nháy (Convergence $N\to\infty$? / Architecture bounds? / Param-Gen in $\infty$-dim?).
& \textit{``Vấn đề mở thứ ba: lý thuyết nền tảng...''}
& Knowledge graph: DL nodes glow + \manim{Line} connect. NO side: fog \manim{FadeIn} + beacon \manim{Flash} tuần tự. \\ \hline
\textbf{[26:05--26:40]} \newline
\textbf{Beat 2 -- Uncertainty Quantification:} Clear beat 1. Hiện function plot: prediction curve (xanh) + confidence band (xanh nhạt mờ). So sánh 2 phương pháp: Monte Carlo (50 đường rải, đỏ, expensive label) vs Conformal/GANO (1 band mượt mà, xanh lá, guarantee label). \manim{SurroundingRectangle} highlight ``Recommended'' cho conformal.
& \textit{``Vấn đề mở thứ tư: lượng hóa bất định...''}
& Prediction \manim{Create} + band \manim{FadeIn}. MC curves \manim{Create} tuần tự. Morph thành conformal band. Highlight box \manim{Create}. \\ \hline
\end{tabularx}

\subsection{Scene 6.3: Active Learning \& RL trong không gian hàm}
\textbf{Timestamp tổng:} \texttt{26:40 -- 27:55}

\setlength{\tabcolsep}{4pt}
\renewcommand{\arraystretch}{1.2}
\begin{tabularx}{\textwidth}{|p{3.6cm}|>{\RaggedRight\arraybackslash}X|p{4.6cm}|}
\hline
\textbf{Visual (Manim)} & \textbf{Audio / Voiceover} & \textbf{Ghi chú Animation} \\ \hline
\textbf{[26:40--27:15]} \newline
\textbf{Beat 1 -- Active Learning:} Hiện ``Active Learning trong không gian hàm''. Trên: Function space với candidate points (scattered). Giữa: Budget bar giảm dần (mỗi sample = cost). Dưới: So sánh Random (100 samples, xám mờ) vs Active (20 samples, xanh sáng). Hiệu ứng: Active samples được ``chọn'' thông minh ở vùng uncertainty cao.
& \textit{``Vấn đề mở thứ năm: active learning trong không gian hàm...''}
& Function space + candidate \manim{Dot} scatter. Budget bar \manim{GrowFromEdge} + animate giảm. So sánh side-by-side: Random grid vs Active targeted. \\ \hline
\textbf{[27:15--27:55]} \newline
\textbf{Beat 2 -- RL trong không gian hàm:} Trên: Classic RL loop (4 ô: State $\mathbf{s}\in\mathbb{R}^n$ $\to$ Policy $\pi$ $\to$ Action $\mathbf{a}\in\mathbb{R}^m$ $\to$ Environment). Chữ trong ô, ký hiệu MathTex ($\pi$, $\Pi$). Mũi tên cong \manim{CurvedArrow} nối Environment $\to$ State. Dấu X đỏ bự giữa mũi tên cong + chữ ``Không áp dụng cho function space!'' (2 dòng). Dưới: Function-Space RL loop (State $u(x)\in\mathcal{U}$ $\to$ Operator Policy $\Pi$ $\to$ Action $a(x)\in\mathcal{A}$ $\to$ PDE Environment). Viền xanh lá sáng, \manim{SurroundingRectangle}. Bullet points: Điều khiển dòng chảy, Can thiệp khí hậu, Tối ưu thiết kế.
& \textit{``Vấn đề mở thứ sáu: reinforcement learning trong không gian hàm...''}
& Classic RL \manim{FadeIn}. \manim{CurvedArrow} angle=$-75^\circ$. Cross \manim{Create} at arc midpoint. Function RL \manim{FadeIn} with glow. Apps \manim{LaggedStart}. \\ \hline
\end{tabularx}

\subsection{Scene 6.4: Thách thức hợp tác \& Recap 3 nguyên tắc}
\textbf{Timestamp tổng:} \texttt{27:55 -- 28:55}

\setlength{\tabcolsep}{4pt}
\renewcommand{\arraystretch}{1.2}
\begin{tabularx}{\textwidth}{|p{3.6cm}|>{\RaggedRight\arraybackslash}X|p{4.6cm}|}
\hline
\textbf{Visual (Manim)} & \textbf{Audio / Voiceover} & \textbf{Ghi chú Animation} \\ \hline
\textbf{[27:55--28:25]} \newline
\textbf{Beat 1 -- Building the Bridge:} \textit{(Đã redesign: thay ``hai nhóm người'' bằng hai thế giới trừu tượng.)} Trái: ``PHYSICS WORLD'' (ô xanh dương, chứa $\nabla\cdot\mathbf{u}=0$, Finite Element, Conservation Laws, 50+ years of solvers). Phải: ``ML WORLD'' (ô tím, chứa Neural Networks, $\mathcal{L}(\theta)$, Gradient Descent, Loss Functions). Giữa: ``THE GAP'' (ô đỏ nhạt, chứa Different Languages, Biased Data, Different Metrics, Skepticism) $\to$ Gap tan biến $\to$ Bridge band xanh lá sáng hiện ra phía dưới chứa ``Function Spaces'', ``Operators'', ``Discretization'', $\mathcal{G}:\mathcal{A}\to\mathcal{U}$. Mũi tên từ 2 thế giới xuống bridge.
& \textit{``Và cuối cùng, một thách thức thực tế: hợp tác với domain experts...''}
& Physics box \manim{Create} + equations \manim{LaggedStart}. ML box \manim{Create} + tools \manim{LaggedStart}. Gap \manim{FadeIn} $\to$ \manim{FadeOut}. Bridge \manim{Create} + \manim{Flash}. Arrows \manim{GrowArrow}. \\ \hline
\textbf{[28:25--28:55]} \newline
\textbf{Beat 2 -- Three Pillars Recap:} Clear beat 1. Hiện 3 ô dọc: (1) Discretization Invariant (icon lưới thô $\to$ mịn phóng to trong ô, xanh dương, tick $\checkmark$). (2) Function Output (đường cong tím + query dots xanh + $\nabla u, \int u\,dx$, tick $\checkmark$). (3) Speed Orders of Magnitude (icon tia sét vàng + 3 con số: 45,000$\times$ Weather / 140,000$\times$ CFD / 700,000$\times$ Carbon, tick $\checkmark$, chữ xanh lá trên 1 dòng). Dưới: impact statement ``Không chỉ là con số --- mở ra những phân tích trước đây không thể làm được''.
& \textit{``Nhìn toàn cảnh lại...''}
& 3 pillar boxes \manim{Create} lần lượt. Icons/motifs bên trong. Tick $\checkmark$ \manim{Write}. Impact text \manim{FadeIn} + \manim{Flash}. Clean \manim{FadeOut}. \\ \hline
\end{tabularx}

\subsection{Scene 6.5: Resources \& Closing remarks}
\textbf{Timestamp tổng:} \texttt{28:55 -- 29:25}

\setlength{\tabcolsep}{4pt}
\renewcommand{\arraystretch}{1.2}
\begin{tabularx}{\textwidth}{|p{3.6cm}|>{\RaggedRight\arraybackslash}X|p{4.6cm}|}
\hline
\textbf{Visual (Manim)} & \textbf{Audio / Voiceover} & \textbf{Ghi chú Animation} \\ \hline
\textbf{[28:55--29:10]} \newline
\textbf{Beat 1 -- The Starter Kit:} Trái (Theory): 2 paper nền tảng Kovachki 2021 và Li 2021 dưới dạng thẻ document. Phải (Code): Cửa sổ Terminal (macOS style) gõ lệnh \texttt{pip install} và \texttt{import FNO} trực quan. Dưới (Hub): Banner CTA màu vàng nổi bật chỉ thẳng người xem xuống phần Mô tả (Description) để lấy toàn bộ tài nguyên.
& \textit{``Nếu bạn muốn bắt đầu khám phá lĩnh vực này...''}
& Theory \manim{FadeIn} with shift. Terminal \manim{FadeIn} + code \manim{Write} inside. Hub banner \manim{FadeIn} + nháy thu hút. Khác với bản cũ, bỏ hoàn toàn QR Code rườm rà. \\ \hline
\textbf{[29:10--29:25]} \newline
\textbf{Beat 2 -- The Final Synthesis:} Clear beat 1. Hiện 3 phép đồng cấu toán học: $\sum \to \int$, $\mathbb{R}^n \to \mathcal{A}$, $W \to \kappa$ và giữ nguyên trên màn hình. Sau đó, Typography focus: Dòng 1 ``Neural Operators không phải thứ gì bí ẩn.'' (trắng). Dòng 2 ``Chúng là sự tổng quát hóa tự nhiên của'' (xám). Dòng 3: 3 keywords highlighted xanh lá ``tổng Riemann | tích phân | và MLP''. Dòng 4 ``lên không gian vô hạn chiều.'' (xám). Fade to black.
& \textit{``Hy vọng video này giúp bạn thấy Neural Operators không phải một thứ gì bí ẩn...''}
& Isomorphisms \manim{Write}. Typography lines \manim{FadeIn}. Keywords \manim{LaggedStart}. Thank you text \manim{FadeIn}. All \manim{FadeOut} to black. \\ \hline
\end{tabularx}

